\section{What \texorpdfstring{$\kappa$}{kappa}-linear symmetry hides: the \texorpdfstring{$\mathbb{F}_p$}{Fp}-shadow of the form}
\label{sec:fp-symmetry}

\paragraph{Recollection.} For $\beta\in\mathrm{Adm}(\omega)$ and phase datum
$\psi$, the Weyl frame is
$\mathcal{F} = \{\mathbb{C}^\times W_\beta(v) : v\in V(\kappa)\} \subset A_{\psi,\beta}(V)$,
and $\mathrm{Aut}_{\mathcal{F}}(A)$ is the group of $\mathbb{C}$-algebra
automorphisms of $A_{\psi,\beta}(V)$ preserving $\mathcal{F}$
(Definition~\ref{def:weyl-frame}); $\mathrm{Aut}_{\mathcal{F}}^\kappa(A) \subseteq \mathrm{Aut}_{\mathcal{F}}(A)$
is the further subgroup whose induced permutation of $V(\kappa)$ is
\emph{$\kappa$-linear} — data \emph{imposed} on top of
$\mathrm{Aut}_{\mathcal{F}}(A)$, not derived from it
(Definition~\ref{def:weyl-frame}'s scope). By
Theorem~\ref{thm:wh-weil-a}, $\mathrm{Aut}_{\mathcal{F}}^\kappa(A)$ surjects
onto $SL_2(\kappa)$ with kernel $\hat V$, at every characteristic.
\provenance{D4, D7, WH-WEIL-a}{---}{---}{---}

This section makes good on the promissory note left in
Definition~\ref{def:weyl-frame}: exactly how much symmetry the imposed
$\kappa$-linearity clause discards. The answer is that $A_{\psi,\beta}(V)$
and its frame $\mathcal{F}$ do not, in fact, remember the $\kappa$-module
structure of $V(\kappa)$ at all — they are built from the bare abelian group
$(V(\kappa),+)$ and the cocycle $\psi\circ\beta$ — so the \emph{full} frame
symmetry group is larger than $SL_2(\kappa)$ whenever $\kappa$ is a proper
extension of its prime field, and the discarded symmetry is measured
exactly.

\statusProved
\begin{theorem}[The twisted algebra does not remember the $\kappa$-structure
of $V(\kappa)$]
\label{thm:wh-symm}
\begin{enumerate}[label=(\alph*)]
\item $A_{\psi,\beta}(V)$ and its frame $\mathcal{F}$ are determined by the
abelian group $(V(\kappa),+)$ together with the $\mathbb{C}^\times$-valued
function $\psi\circ\beta$ on $V(\kappa)\times V(\kappa)$ alone; nothing in
either object refers to scalar multiplication by $\kappa$;
\item consequently, dropping the imposed $\kappa$-linearity clause of
Definition~\ref{def:weyl-frame}, the frame automorphisms with merely
$\mathbb{F}_p$-linear induced permutation $g$ of $V(\kappa)$ surject, with
the \emph{same} kernel $\hat V$, onto the isometry group
$Sp(V(\kappa),\psi\circ\omega) \cong Sp_{2m}(\mathbb{F}_p)$ of the
$\mathbb{F}_p$-valued alternating form $\psi\circ\omega$; so
$|\mathrm{Aut}_{\mathcal{F}}(A)| = q^2\cdot|Sp_{2m}(\mathbb{F}_p)|$;
\item $SL_2(\kappa) \subseteq Sp_{2m}(\mathbb{F}_p)$, with index exactly $1$
when $m=1$ and strictly greater than $1$ as soon as $m>1$: by complete
enumeration, $|SL_2(\kappa)| = 6,24,60,120,504,720$ against
$|Sp_{2m}(\mathbb{F}_p)| = 6,24,720,120,1451520,51840$ at
$q=2,3,4,5,8,9$ respectively — indices $1,1,12,1,2880,72$. So
Definition~\ref{def:weyl-frame}'s $\kappa$-linearity clause discards genuine
symmetry the algebra has, whenever $\kappa$ is a proper extension of its
prime field.
\end{enumerate}
\end{theorem}

\begin{scope}
This theorem does not claim $Sp_{2m}(\mathbb{F}_p)$ acts by automorphisms
preserving anything beyond the frame $\mathcal{F}$ and the abelian group
structure — in particular it says nothing about whether this larger
symmetry is compatible with the $\kappa$-linear structure used elsewhere in
this labbook (Theorem~\ref{thm:wh-form}(f) shows, to the contrary, that the
\emph{$\kappa$-valued} form $\omega$ itself sees only $SL_2(\kappa)$ among
$\mathbb{F}_p$-linear maps: the enlargement here is possible only because
$\mathrm{Aut}_{\mathcal{F}}(A)$ is built from the strictly weaker
$\mathbb{F}_p$-valued shadow $\psi\circ\omega$). At the prime field, $m=1$,
the two groups coincide and nothing is discarded; this is the campaign's
guiding path.
\end{scope}

\provenance{D1, D7, WH-WEIL-a, WH-FORM, WH-SYMM}{theory/wh-kappa-choice.md \S9}{theory/checks/fp\_symmetry.py (S1,S2,S3)}{theory/verdicts/wh-kappa-adjudication.md}

\begin{proof}
(a) By Definition~\ref{def:weyl-algebra}, the structure constants of
$A_{\psi,\beta}(V)$ in the basis $\{W_\beta(v)\}$ are exactly the values
$\psi(\beta(v,v'))$, indexed by pairs of elements of the abelian group
$(V(\kappa),+)$; by Definition~\ref{def:weyl-frame}, $\mathcal{F}$ is the
set of $\mathbb{C}^\times$-lines through this same basis, again indexed only
by $(V(\kappa),+)$. Neither construction refers to multiplication by
elements of $\kappa$.

(b) Repeat the argument of Theorem~\ref{thm:wh-weil-a}'s proof
\emph{without} imposing $\kappa$-linearity on $g$: frame preservation still
gives $\alpha(W_\beta(v))=\lambda(v)W_\beta(gv)$ with $g$ an (a priori only
additive) bijection of $V(\kappa)$ and $\lambda$ satisfying relation $(*)$
there, since that derivation used only additivity of $g$, never
$\kappa$-linearity. Applying $\alpha$ to Theorem~\ref{thm:wh-comm}(c), an
identity entirely in $\psi\circ\omega$, gives
$\psi\big(\omega(gv,gv')\big)=\psi\big(\omega(v,v')\big)$ for all $v,v'$ —
i.e.\ $g\in Sp(V(\kappa),\psi\circ\omega):=\{g \text{ additive} : \psi\circ\omega\circ(g\times g)=\psi\circ\omega\}$
— rather than the stronger $\omega\circ g=\omega$ used before. The
surjectivity argument via Fact~\ref{fact:triv} applies verbatim to any such
$g$ (it never used $\kappa$-linearity, only that $(V(\kappa),+)$ has
exponent $p$), so every $g\in Sp(V(\kappa),\psi\circ\omega)$ lifts; the
kernel computation ($g=\mathrm{id}$ forces $\lambda\in\hat V$, identified
with $V(\kappa)$ via $u\mapsto\psi(\omega(u,\cdot))$) is unchanged, giving
the same kernel $\hat V$ of order $q^2$. Nondegeneracy of $\psi\circ\omega$
as an $\mathbb{F}_p$-bilinear form on the $2m$-dimensional $\mathbb{F}_p$-space
$V(\kappa)$: if $\psi(\omega(v,\cdot))\equiv 1$ then $\omega(v,\cdot)\equiv0$
(Fact~\ref{fact:val}, since $\omega(v,\cdot)$ is either $0$ or onto $\kappa$,
being $\kappa$-linear, and $\kappa\not\subseteq\ker\psi$), so $v=0$ by
Theorem~\ref{thm:wh-form}(c). A nondegenerate alternating form on a
$2m$-dimensional space over $\mathbb{F}_p$ has isometry group (by
definition, matching the notation \texttt{notation.md} fixes for
$Sp_{2m}(\mathbb{F}_p)$) exactly $Sp_{2m}(\mathbb{F}_p)$.

(c) If $g\in SL_2(\kappa)$ then $\omega\circ(g\times g)=\omega$
(Theorem~\ref{thm:wh-form}(e)), so certainly
$\psi\circ\omega\circ(g\times g)=\psi\circ\omega$, i.e.\
$SL_2(\kappa)\subseteq Sp(V(\kappa),\psi\circ\omega)$. This containment is
exactly the gap between Theorem~\ref{thm:wh-form}(f) — which shows the
$\mathbb{F}_p$-linear isometries of the \emph{$\kappa$-valued} $\omega$ are
no larger than $SL_2(\kappa)$ — and the present, strictly weaker condition
of isometry for the $\mathbb{F}_p$-valued shadow $\psi\circ\omega$ alone: an
$\mathbb{F}_p$-linear $g$ can preserve $\psi\circ\omega$ without preserving
$\omega$ itself. Complete enumeration
(\texttt{theory\slash{}checks\slash{}fp\_symmetry.py}, gates S1 and S2) computes both
orders directly at $q=2,3,4,5,8,9$, matching the closed form
$|Sp_{2m}(\mathbb{F}_p)| = p^{m^2}\prod_{i=1}^m(p^{2i}-1)$ in every case,
and confirms the indices $1,1,12,1,2880,72$; at $m=1$ (the prime field) the
two group orders coincide, since $Sp_2(\mathbb{F}_p)=SL_2(\mathbb{F}_p)$ by
definition in dimension $2$, and nothing is discarded.
\end{proof}

\begin{remark}[The characteristic-two orthogonal phenomenon survives the
enlargement]
The $\mu_2$-phase criterion of Theorem~\ref{thm:wh-weil-c}
($g\in O(Q_\beta)$) makes sense verbatim for merely $\mathbb{F}_p$-linear
$g$, at $p=2$: $\psi\circ Q_\beta\circ g=\psi\circ Q_\beta$ is a condition
on an $\mathbb{F}_p$-valued quadratic form $\psi\circ Q_\beta$ with polar
form $\psi\circ\omega$, so the relevant subgroup inside the enlarged
symmetry group of Theorem~\ref{thm:wh-symm} is
$O(\psi\circ Q_\beta) \cap Sp_{2m}(\mathbb{F}_p)$. Direct enumeration
(gate S3) finds its order to be $2,72,40320$ at $q=2,4,8$, of index
$3,10,36$ in $Sp_{2m}(\mathbb{F}_p)$ — exactly the \emph{same} index that
$|O(Q_\beta)\cap SL_2(\kappa)|=2,6,14$ has inside $SL_2(\kappa)$
(Theorem~\ref{thm:wh-weil-c}). So the characteristic-two orthogonal
phenomenon of \S\ref{sec:weil} is not an artefact of the imposed
$\kappa$-linearity in Definition~\ref{def:weyl-frame}: it survives, at the
same relative index, inside the larger, intrinsic symmetry group
$Sp_{2m}(\mathbb{F}_p)$ that Theorem~\ref{thm:wh-symm} exhibits. This
observation is recorded as part of the same source theorem as
Theorem~\ref{thm:wh-symm}, under the same provenance
(\texttt{theory\slash{}wh-kappa-choice.md} \S9, step $\langle1\rangle4$).
\end{remark}

\begin{remark}[Reading]
Two independent non-canonicities are now on record, and they are logically
independent of each other. Theorem~\ref{thm:wh-beta-type} of
\S\ref{sec:polarizing-cocycle} shows the system does not remember which
polarizing cocycle $\beta$ was chosen, in a way that is invisible at odd
characteristic and genuinely two-valued at $p=2$. This theorem shows the
system does not remember the $\kappa$-module structure of $V(\kappa)$
itself, at every characteristic and every $m>1$, regardless of $\beta$.
Both are statements about what the definition actually depends on, in the
sense of \texttt{PRD.md}'s ``canonical'' requirement: any assignment
$\mathrm{Spec}\,R \mapsto$ quantum system general enough to restrict to this
one on $\mathrm{Spec}\,\mathbb{F}_q$ carries both phenomena, not as defects,
but as exactly the symmetry and the choice the construction is built from.
\end{remark}
